% The body of the merged paper.  Every number in this file is checked by scripts/j_checknums.py
% against results/json; the two data tables are generated by paper/gen_tables.py.
\section{Introduction}
\label{sec:intro}

For a graph $G$ the Shannon capacity is $\Theta(G)=\sup_d \alpha(G^{\boxt d})^{1/d}$
\cite{shannon1956}. For the seven-cycle it is unknown. The Lov\'asz bound gives
$\Theta(C_7)\le\vartheta(C_7)=3.3176672073940953927\ldots$ \cite{lovasz1979}, and the lower
bound moved only in its fourth decimal place between 2019 and July 2026 (Table~\ref{tab:chain}).

\begin{table}[t]
\centering\footnotesize\setlength{\tabcolsep}{3pt}
\begin{tabular}{@{}llll@{}}
\toprule
date & $\Theta(C_7)\ge$ & from & witness \\
\midrule
1971--2015 & $343^{1/5}$, then $350^{1/5}$ & \cite{baumert1971,mathew2017} & linear code; search \\
2019 & $3.2578659667835159504124\ldots$ & \cite{polak2019} & $\alpha(C_7^{\boxt 5})\ge 367$ \\
Jul 2026 & $3.2580207372932453595278\ldots$ & \cite{itty2026} & $\alpha(C_7^{\boxt 10})\ge 134753$ \\
Jul 2026 & $3.2587891539086910161967650155206769\ldots$ & \cite{gao2026} & gadget product \\
Jul 2026 & $3.2588053698854655725829750306238067\ldots$ & \cite{bpz2026} & stronger base gadget \\
Aug 2026 & $3.2588326203532663091215390518104754\ldots$ & \cite{tandon2026} & heterogeneous recursion \\
\bottomrule
\end{tabular}
\caption{Lower bounds on $\Theta(C_7)$. The last four entries start from the same $367$-word
code and the same eight private pairs, and differ in the recursion built on top.}
\label{tab:chain}
\end{table}

The structure behind the last four rows is Gao's \emph{gadget} \cite{gao2026}: an independent set
(the code), a family of \emph{private pairs} --- a vertex $q$ outside the code with a single
neighbour $r$ in it --- two transversals of their endpoints, and an auxiliary independent set.
A product lemma multiplies gadgets, and the bound it yields is most sensitive to the number $t$
of private pairs. Evaluated at hypothetical base profiles $(a,t,s,o)$, the homogeneous recursion
of \cite{gao2026,bpz2026} gives
\begin{align*}
(367,8,367,322)&\mapsto 3.2588053698854655725829750306238067\ldots,\\
(367,9,367,321)&\mapsto 3.2590062007738521305802\ldots,
\end{align*}
so one further private pair would be worth $1.736\cdot 10^{-4}$ over the current record, while
the record itself \cite{tandon2026} exceeds what the homogeneous recursion extracts from the
\emph{same} base gadget by only $2.725\cdot10^{-5}$. A ninth pair would therefore improve the
state of the art under any of the published recursions, by several times the spread between
them.

This paper studies the number of private pairs as a characteristic of a code, and explains why
that number is not, by itself, a lever. It does so in three steps, from the general to the
particular. First (Section~\ref{sec:cost}), the cost a pair imposes on the auxiliary set is
conserved: the forbidden and the penalised regions a pair creates are the intersection and the
union of two closed neighbourhoods, so their sizes add up to $2\cdot 3^d$, and the shape of a
pair can move cost between them but not lower it. Second (Section~\ref{sec:pscode}), for the
Polak--Schrijver code itself the count is at its stop: it admits exactly eight private pairs, and
we prove that no ninth exists. Third (Sections~\ref{sec:count}--\ref{sec:gap}), for the codes
with nine and ten pairs that do exist, the forbidden regions are fixed by the geometry of the pairs
rather than by the code, and against every one of them the best auxiliary set we could build is
one word short.

The status of each claim is stated where it is made, and three kinds are kept apart: \emph{proved}
(by hand, with any finite check it rests on named); \emph{exhaustive computation} (a complete
enumeration of a finite object, reproducible from the published sets); and \emph{search result}
(the best found, which is not a bound). The results here were first deposited as two preprints
\cite{oktiabrev2026a,oktiabrev2026b}\footnote{This paper merges and reorganises the preprints
\cite{oktiabrev2026a} (the theorem of Section~\ref{sec:pscode} and the computations around it) and
\cite{oktiabrev2026b} (Sections~\ref{sec:cost}, \ref{sec:shapes}, \ref{sec:regions} and
\ref{sec:gap}).}.

\section{Gadgets, private pairs and the two regions}
\label{sec:defs}

Vertices of $C_n^{\boxt d}$ are the elements of $\mathbb{Z}_n^d$; two distinct vertices are
\emph{confusable} iff their circular distance is at most $1$ in every coordinate. $N(S)$ is the
\emph{closed} neighbourhood of a set $S$, as in \cite{gao2026,tandon2026}; for $n\ge3$ and a
single vertex, $|N(v)|=3^d$. A set is independent if no two of its elements are confusable.

\begin{definition}[private pair \cite{gao2026}]
\label{def:pair}
Let $I$ be independent. A pair $(r,q)$ with $r\in I$, $q\notin I$ is a \emph{private pair} for
$I$ if $N(q)\cap I=\{r\}$; $r$ is its \emph{centre} and $q$ its \emph{private neighbour}. The
\emph{candidates} of $I$ are all such pairs.
\end{definition}

\begin{definition}[gadget \cite{gao2026}]
\label{def:gadget}
A gadget consists of an independent set $I$ of size $a$ (the \emph{code}); private pairs
$(r_i,q_i)$, $1\le i\le t$, with pairwise disjoint endpoint sets; two complementary transversals
$\PH,\PV$ of the endpoints --- for every $i$ exactly one of $r_i,q_i$ lies in $\PH$ and the other in
$\PV$ --- both independent; and an independent auxiliary set $X$ with
$X\cap N(\PH)\cap N(\PV)=\emptyset$. With $O=X\setminus(N(\PH)\cup N(\PV))$, $H=X\cap N(\PH)$,
$V=X\cap N(\PV)$, the profile is $(a,t,s,o,h,v)=(a,t,|X|,|O|,|H|,|V|)$.
\end{definition}

We name the two sets in this definition:
\begin{align*}
  F&=N(\PH)\cap N(\PV) && \text{(the forbidden region)},\\
  U&=N(\PH)\cup N(\PV) && \text{(the penalised region)}.
\end{align*}
$X$ must avoid $F$, and every word of $X$ in $U$ is lost from $o$. Write $t^*(I)$ for the
largest $t$ for which a gadget with code $I$ exists. The \emph{shape} of a pair is
$\delta=q-r\in\{0,\pm1\}^d\setminus\{0\}$, and $k$ is its number of non-zero coordinates.

The published base gadget has $I=I_0$, the $367$-word set printed in the appendix of
\cite{polak2019}, and $t=8$ with the eight pairs of \cite{itty2026}; its profile is
$(367,8,367,321,26,20)$ in \cite{gao2026} and $(367,8,367,322,26,19)$ in \cite{bpz2026,tandon2026}.

\section{Cost is conserved}
\label{sec:cost}

\begin{lemma}[proved]
\label{lem:conserve}
Let $n\ge3$ and let $\PH,\PV$ be any two sets of vertices of $C_n^{\boxt d}$. Then
\[
  |F|+|U| \;=\; |N(\PH)|+|N(\PV)| \;\le\; 3^d\,(|\PH|+|\PV|),
\]
with equality iff the closed neighbourhoods of the vertices inside each of $\PH$ and $\PV$ are
pairwise disjoint. In particular, for a single pair, $\PH=\{r\}$ and $\PV=\{q\}$,
\[
  |N(r)\cap N(q)|+|N(r)\cup N(q)| \;=\; 2\cdot 3^d .
\]
\end{lemma}

\begin{proof}
$|A\cap B|+|A\cup B|=|A|+|B|$ for finite sets, with $A=N(\PH)$, $B=N(\PV)$; and
$|N(S)|\le\sum_{s\in S}|N(s)|=3^d|S|$, since for $n\ge3$ the three values $a-1,a,a+1$ of every
coordinate are distinct.
\end{proof}

The lemma is inclusion--exclusion and holds for any two vertices, private or not; what matters is
its consequence for gadgets. A pair contributes $N(r)\cap N(q)$ to $F$ and $N(r)\cup N(q)$ to $U$,
and the two add up to $2\cdot 3^d$, which is $486$ in $C_7^{\boxt 5}$. \emph{The shape of a pair
redistributes its cost between the forbidden and the penalised region; it does not lower it.} A
shape that makes $F$ small makes $U$ large by the same amount, and $U$ is what $o$ is paid from.

\begin{lemma}[proved]
\label{lem:shape}
For $n\ge4$ and a pair of shape $k$, $|N(r)\cap N(q)|=3^{d-k}\,2^{k}$.
\end{lemma}

\begin{proof}
Coordinatewise, $\{a-1,a,a+1\}\cap\{a+\delta-1,a+\delta,a+\delta+1\}$ has $3$ elements if
$\delta=0$ and $2$ if $\delta=\pm1$, for $n\ge4$; closed neighbourhoods in the strong product are
products of the coordinate neighbourhoods.
\end{proof}

Both lemmas were also checked by direct count for every shape with $d=3,4,5$ and $n=3,4,5,7,9$
($n=3$ satisfies Lemma~\ref{lem:conserve} and, as it must, not Lemma~\ref{lem:shape}). For
families the inequality is strict in practice: the nine-pair configurations of
Section~\ref{sec:regions} have $|F|+|U|$ between $4050$ and $4080$, against
$2\cdot9\cdot3^5=4374$.

\section{The five shapes}
\label{sec:shapes}

\begin{table}[t]
\centering\small
\begin{tabular}{@{}crrrrr@{}}
\toprule
$k$ & shapes & $|N(r)\cap N(q)|$ & $|N(r)\cup N(q)|$ & covered cells & largest set \\
 & & (forbidden) & (penalised) & in box$(q)$ & found \\
\midrule
1 & 10 & 162 & 324 & 16 & 366 \\
2 & 40 & 108 & 378 & 8 & 367 \\
3 & 80 & 72 & 414 & 4 & 366 \\
4 & 80 & 48 & 438 & 2 & 366 \\
5 & 32 & 32 & 454 & 1 & 366 \\
\bottomrule
\end{tabular}
\caption{The five shapes of a private pair in $C_7^{\boxtimes 5}$, by the number $k$ of non-zero coordinates of $q-r$. The first four columns are counted directly; the last is the largest independent set the exhaustive local construction of Section~\ref{sec:shapes} produced carrying a private pair of that shape.}
\label{tab:shapes}
\end{table}


The $242$ shapes in $C_7^{\boxt 5}$ fall into five orbits under the stabiliser of a vertex in
$(D_7)^5\rtimes S_5$, one per value of $k$ (Table~\ref{tab:shapes}; exhaustive computation). The
cost in $F$ \emph{falls} with $k$, so the two-coordinate shape is in the middle, not at the minimum.
Every shape is realisable as a private pair --- the expensive $k=1$ included: all seven candidates
of the Mathew--\"Osterg{\aa}rd $350$-word set \cite{mathew2017} have $k=1$.

The last column of Table~\ref{tab:shapes} is an exhaustive computation over a finite list of
inputs: for every word $r$ of every code in Table~\ref{tab:codes} and of the swap plateau of
Section~\ref{sec:regions}, and every shape, put $q=r+\delta$, delete the other code words
confusable with $q$, and add back an exact maximum independent set of the freed vertices that keep
$q$ private. In everything this produces, shapes other than $k=2$ occur only one word below $367$.

The column of covered cells explains why. A word $w$ occupies the box $w+\{0,1\}^5$ of $32$
cells, and for $n\ge4$ two vertices are non-confusable exactly when their boxes are disjoint; a
private $q$ meets the code's boxes only inside box$(r)$, in $2^{5-k}$ cells. A pair of shape
$k\ge3$ therefore needs a box with at most $4$ covered cells. In every $367$-word code we can
reach, every box outside the code has at least $5$.

\section{The Polak--Schrijver code has exactly eight private pairs}
\label{sec:pscode}

The definition of a gadget quantifies over transversals, which makes $t^*$ look hard to bound.
It is not: the transversal condition collapses to a condition on the private neighbours alone.

\begin{lemma}[no mixing; proved]
\label{lem:nomix}
Let $(r_i,q_i)$ and $(r_j,q_j)$ be private pairs for $I$ with $r_i\neq r_j$. Then $r_i$ is
confusable with neither $r_j$ nor $q_j$.
\end{lemma}

\begin{proof}
$r_i,r_j\in I$ and $I$ is independent, so $r_i\not\sim r_j$. If $r_i\sim q_j$ then
$r_i\in N(q_j)\cap I=\{r_j\}$, so $r_i=r_j$, contradicting the hypothesis.
\end{proof}

\begin{corollary}[proved]
\label{cor:bip}
Let $(r_i,q_i)$, $1\le i\le t$, be private pairs for $I$ with pairwise distinct centres, and let
$\sigma:\{1,\dots,t\}\to\{\mathrm H,\mathrm V\}$ assign to each pair the transversal holding $q_i$.
Then the endpoint sets are pairwise disjoint, and $\PH,\PV$ are both independent if and only if
$\sigma$ is a proper $2$-colouring of the graph $\Gamma$ on $\{q_1,\dots,q_t\}$ whose edges are the
confusable pairs $q_i\sim q_j$. Consequently $t^*(I)$ is the largest number of private pairs with
distinct centres whose $q$-graph $\Gamma$ is bipartite.
\end{corollary}

\begin{proof}
Distinct centres give $r_i\neq r_j$; $q_i\neq q_j$ because $N(q_i)\cap I=\{r_i\}$ pins the centre;
and $r_i\neq q_j$ because $r_i\in I$, $q_j\notin I$. So the endpoint sets are disjoint. Two $r$'s are
never confusable, and an $r$ and a $q$ of different pairs are never confusable by
Lemma~\ref{lem:nomix}; $r_i$ and $q_i$ are confusable but lie in different transversals by
construction. Hence a transversal can fail to be independent only by containing two confusable
private neighbours, which is a monochromatic edge of $\Gamma$.
\end{proof}

So the only objects to examine are the candidates of $I$, a finite list that can be written down.

\begin{proposition}[exhaustive computation]
\label{prop:count}
Let $I_0\subseteq\Zs$ be the $367$-word Polak--Schrijver code as printed in \cite{polak2019}. Of the
$16807-367=16440$ vertices outside $I_0$, exactly $8$ have one neighbour in $I_0$, $254$ have two,
and $16178$ have three or more; none has zero. The eight are exactly the private neighbours
$q_0,\dots,q_7$ of \cite{itty2026,gao2026}, with the centres those papers give them.
\end{proposition}

\begin{proof}
One pass over $\Zs$, counting for each vertex outside $I_0$ its neighbours in $I_0$, with the integer
predicate of Section~\ref{sec:defs}. The empty zero-bucket says that $I_0$ is a maximal independent
set, which is independently known \cite{polak2019}. The identification with the published pairs is a
set comparison.
\end{proof}

\begin{theorem}[proved]
\label{thm:main}
$t^*(I_0)=8$. Moreover the admissible transversal assignments are exactly the proper
$2$-colourings of a graph with five connected components, $16$ of them up to exchanging $\PH$ and
$\PV$, and one of them is the assignment used in \cite{itty2026,gao2026,bpz2026,tandon2026}.
\end{theorem}

\begin{proof}
By Corollary~\ref{cor:bip}, $t^*(I_0)$ is computed from the candidate list of
Proposition~\ref{prop:count}, which has eight entries with eight distinct centres; so
$t^*(I_0)\le 8$. Its graph $\Gamma$ has, in the numbering of \cite{gao2026}, exactly the three edges
\[
  q_0\sim q_1,\qquad q_2\sim q_6,\qquad q_3\sim q_5,
\]
a matching on six of the eight vertices with two isolated vertices. A matching is bipartite, so all
eight pairs may be used at once and $t^*(I_0)=8$. $\Gamma$ has $3+2=5$ components, hence $2^5$ proper
$2$-colourings, $2^4=16$ up to a global swap; that the published assignment is one of them is
checked directly.
\end{proof}

\begin{corollary}
No gadget whose code is $I_0$ has $t\ge 9$. In particular the profile $(367,9,367,\cdot)$ is not
reachable from the Polak--Schrijver code, under the homogeneous recursion of
\cite{gao2026,bpz2026} or any recursion that takes a gadget over this code as its base, including
\cite{tandon2026}.
\end{corollary}

Lemma~\ref{lem:nomix} and Corollary~\ref{cor:bip} are proved by hand; the two finite facts they are
applied to --- the candidate list and the three edges --- are computations over $16807$ vertices
that any reader can repeat from the printed set.

Two further questions around the theorem are settled by exhaustive computation. \emph{All codes the
Polak--Schrijver construction can produce:} step (v) of \cite{polak2019} extends a core of $327$
words by a maximum independent set of an extension graph with $71$ vertices and $85$ edges, whose
independence number is $40$; that maximum is not unique, and enumerating all of them gives exactly $8$
distinct codes of size $367$, with candidate counts
\[
  5,\;6,\;6,\;6,\;7,\;7,\;7,\;8 ,
\]
the unique code attaining $8$ being the printed one. \emph{Codes of other sizes:} a gadget's code need
not be maximum, but matching the bound of $(367,8,367,321)$ through the published recursion would
require $13$ private pairs at $a=366$, $16$ at $a=365$ and $23$ at $a=360$.

\section{The count is a property of the code}
\label{sec:count}

\begin{table}[t]
\centering\footnotesize\setlength{\tabcolsep}{3pt}
\begin{tabular}{@{}lrrrrrl@{}}
\toprule
code & size & cand. & $t^*$ & colourings & $|F|$ at $t^*$ & by $k=1,\dots,5$ \\
\midrule
linear code \cite{baumert1971} & 343 & 0 & 0 & 0 & -- & 0, 0, 0, 0, 0 \\
Mathew--\"Osterg{\aa}rd \cite{mathew2017} & 350 & 7 & 7 & 64 & 1022--1118 & 7, 0, 0, 0, 0 \\
auxiliary-set search, $t=9$ & 363 & 20 & 18 & 96 & 2171--2308 & 2, 16, 2, 0, 0 \\
repaired auxiliary set, $t=10$ & 365 & 13 & 13 & 128 & 1582--1726 & 0, 13, 0, 0, 0 \\
repaired auxiliary set, $t=9$ & 366 & 12 & 12 & 64 & 1460--1561 & 0, 12, 0, 0, 0 \\
pipeline code 4 & 367 & 5 & 5 & 16 & 552--568 & 0, 5, 0, 0, 0 \\
$I_0$, our rebuild & 367 & 6 & 6 & 16 & 686--708 & 0, 6, 0, 0, 0 \\
pipeline code 0 & 367 & 6 & 6 & 16 & 686--708 & 0, 6, 0, 0, 0 \\
pipeline code 5 & 367 & 6 & 6 & 16 & 686--708 & 0, 6, 0, 0, 0 \\
pipeline code 6 & 367 & 6 & 6 & 16 & 688--708 & 0, 6, 0, 0, 0 \\
pipeline code 1 & 367 & 7 & 7 & 16 & 820--844 & 0, 7, 0, 0, 0 \\
pipeline code 2 & 367 & 7 & 7 & 16 & 822--848 & 0, 7, 0, 0, 0 \\
pipeline code 7 & 367 & 7 & 7 & 16 & 822--848 & 0, 7, 0, 0, 0 \\
pipeline, best & 367 & 8 & 8 & 16 & 956--984 & 0, 8, 0, 0, 0 \\
$I_0$, as printed \cite{polak2019} & 367 & 8 & 8 & 16 & 956--984 & 0, 8, 0, 0, 0 \\
pipeline code 3 & 367 & 8 & 8 & 16 & 956--984 & 0, 8, 0, 0, 0 \\
$X$ of \cite{gao2026}, as a code & 367 & 9 & 9 & 16 & 1098--1128 & 0, 9, 0, 0, 0 \\
ten-candidate code & 367 & 10 & 10 & 16 & 1240--1268 & 0, 10, 0, 0, 0 \\
$o=322$ auxiliary set & 367 & 10 & 10 & 16 & 1240--1268 & 0, 10, 0, 0, 0 \\
\bottomrule
\end{tabular}
\caption{The candidate count of every code we hold in $C_7^{\boxtimes 5}$. Candidates are the vertices with exactly one neighbour in the code; $t^*$ is the largest usable family (Corollary~\ref{cor:bip}); colourings are counted up to exchanging $\PH$ and $\PV$; $|F|$ ranges over them. Exhaustive for each file.}
\label{tab:codes}
\end{table}


Counting the candidates of every code we hold gives Table~\ref{tab:codes}, a statistic that, as far as
we can find, no paper reports. The counts do not follow size: the linear $343$-word code
\cite{baumert1971} has none, the $350$-word set of \cite{mathew2017} has seven, the printed
$367$-word code has eight, and an independent rebuild of it, agreeing with the printed set in $365$
of its words, has six. Among the $367$-word codes the counts run from $5$ to $10$, every candidate has
shape $k=2$, and for a fixed code the forbidden region varies over the admissible colourings by at
most $3.21\%$. No paper we know of reports a code with nine or ten candidates; one of the codes that has nine is
the auxiliary set of the published gadget \cite{gao2026}, read as a code.

\section{Sixteen forbidden regions}
\label{sec:regions}

A private pair is also a move: $I\mapsto I-r+q$ keeps size and independence; it is the exchange used
in \cite{tandon2026}. Closing every $367$-word code we hold under these moves, and under exhaustive
exchanges of two and three words, gives a finite plateau that falls into $51$ classes under
$\mathrm{Aut}(C_7^{\boxt 5})$ (exhaustive computation; exchanges of four words were time-limited
and found nothing new). Each connected component is a product of five independent three-state
switches, $3^5=243$ codes; the five switches sit on the five cyclically adjacent pairs of coordinates
and are permuted by an automorphism of order $5$ that cycles the coordinates. The number of
components depends on the starting codes and is not structural.

Two classes have nine candidates, and one class has ten. An automorphism $g$ moves the forbidden
region to $g\cdot F$ and leaves $\alpha(C_7^{\boxt5}\setminus F)$ unchanged, so the variables are the
code up to $\mathrm{Aut}$ and the colouring; sweeping automorphisms of the code adds nothing. All
$208$ configurations with $t\ge9$ give, up to automorphism, exactly $16$ forbidden regions at $t=9$
and $4$ at $t=10$, with
\[
  1098\le|F|\le1128\ (t=9),\qquad 1240\le|F|\le1268\ (t=10),
\]
and \emph{each of the sixteen is produced by every nine-pair class} (exhaustive computation). The
forbidden region is a function of the geometry of the nine pairs, which are the same switches in
every code, and not of the rest of the code. Of a nine-pair region, $972=9\cdot108$ vertices come
from the pairs themselves and $126$ to $156$ from cross terms between pairs.

\section{One word short}
\label{sec:gap}

\paragraph{At $t=8$.} With $t=8$ forced on $I_0$, the remaining parameter is $X$. The candidates for
$X$ are the eight pipeline codes and their images under $(D_7)^5\rtimes S_5$, a group of
$14^5\cdot 120=64{,}538{,}880$ elements; with every admissible colouring the sweep covers
$8{,}260{,}976{,}640$ triples (source, colouring, automorphism), exhaustively, since for fixed
permutation and reflections all translations are scored at once by a cyclic cross-correlation. The
forbidden intersection is never empty:

\begin{center}\small
\begin{tabular}{@{}lrrrrrrrr@{}}
\toprule
words of $X$ in $F$ & 0 & 1 & 2 & 3 & 4 & 5 & 6 & 7 \\
\midrule
triples & $0$ & $4$ & $33$ & $101$ & $209$ & $358$ & $965$ & $6{,}060$ \\
\bottomrule
\end{tabular}
\end{center}

This explains a step of the published construction stated there without comment: $X$ is the image
$T(I_0)$ under $T(w)=(2-w_1,w_3,w_0,2-w_2,w_4)$ with the single word $(2,4,6,3,5)$ replaced by
$(1,5,6,3,5)$ \cite{gao2026,tandon2026}. The replacement is unavoidable, and $(2,4,6,3,5)$ is exactly
the one word of $T(I_0)$ in $F$. Exact repair of the near misses reaches $s=367$, $o=322$, the profile
of \cite{bpz2026}, recovered from another direction and not improved. A $367$-word admissible
auxiliary set was found under exactly one of the $16$ colourings, the one the literature uses; for the
other fifteen this is a search result ($1500$\,s of fixed-cardinality search each), not a proof.

\paragraph{At $t=9$ and $t=10$.} Against each of the $20$ regions of Section~\ref{sec:regions}, every
image under $(D_7)^5\rtimes S_5$ of one representative of each of the $51$ classes was scored
exhaustively, $65{,}829{,}657{,}600$ images in all. No image avoids any region; the fewest words of an
image in $F$ is $1$. Every image within three words of admissible was repaired exactly, and the best
results are
\[
  366 \text{ words at } t=9,\qquad 365 \text{ at } t=10 .
\]
All $108$ repaired $366$-word sets are exactly $3$-opt optimal inside $G_F=C_7^{\boxt5}\setminus F$: no
window of at most three words can be traded for one more, decided by counting blockers. Only $2$ or
$3$ allowed vertices are blocked by a single word of such a set.

\emph{Status.} This is a computational result, not a proof that $367$ words do not fit outside $F$.
The cheap upper bounds cannot see $F$: the Lov\'asz bound gives $401$ whether $F$ is empty or removes
$1268$ vertices, because it bounds the whole graph, and a box-volume count is weaker still. An honest
bound needs $\vartheta$ of the induced subgraph on $15{,}679$ to $15{,}709$ vertices, and $F$ breaks the
symmetry that would reduce it.

\paragraph{Why a cheaper shape does not obviously help.} With $|U|$ taken without overlaps, the
fraction of its random share that $X$ may touch at $t=9$ and $367$ words is $0.879$, $0.754$, $0.688$,
$0.651$, $0.628$ for $k=1,\dots,5$; the best auxiliary sets on record reach $0.7678$. By
Lemma~\ref{lem:conserve} the shapes that would shrink $F$ are the ones that ask most of $o$, and at
$366$ words no shape allows more than $0.68$ of the random share at $t=9$, $13$ or $16$. These are
distances to the best known search, not impossibility results.

\section{What remains open}
\label{sec:open}

On $I_0$ there is no ninth pair (Theorem~\ref{thm:main}). On the codes with nine and ten pairs that we
hold, the best auxiliary set falls one word short at nine pairs and two at ten, against forbidden
regions fixed by the geometry of the pairs; that is a statement about the codes we hold, not about all
codes. Two things would change it.

\begin{enumerate}
\item \emph{A $367$-word code with a deeper hole}: a box covered in at most $4$ cells, which a pair of
shape $k\ge3$ needs. None of the codes we can reach has one; finding one requires $367$-word codes from
outside the known family, which is the open problem around $\alpha(C_7^{\boxt 5})$ itself, where even
$368$ is open.
\item \emph{An upper bound that sees $F$}, of $\vartheta$ strength on the induced subgraph, which would
turn the one-word gap of Section~\ref{sec:gap} into a theorem or refute it.
\end{enumerate}

\section{Methods and reproducibility}
\label{sec:repro}

All computations are exact integer computations. Independence is checked by a standalone C verifier:
for $n\ge4$ two distinct vertices of $C_n^{\boxt d}$ are adjacent exactly when their boxes meet, so a
set is independent iff its $|I|\cdot 2^d$ box cells are all distinct; for a $367$-word set that is
$11{,}744$ cells. The verifier is held to cases where it can fail: $1177$ single-vertex corruptions
against independently computed ground truth, and a deliberately defective build that must be rejected.
Each further estimator was calibrated in the same way before use: the new group sweep reproduces the
$t=8$ histogram above exactly, the repair reaches $367$ at $t=8$, the window search recovers $367$ from
deliberately damaged sets, and the automorphism checker rejects a non-automorphism. Predictions were
registered before each computation.

The software was written, and the computations run, through Claude Code (Anthropic) under the author's
direction, following preregistered stage plans; every number in this paper is reconciled by script
against the stored results, an unsourced number counting as a failed build. Theorem~\ref{thm:main} and
the lemmas do not rest on the software: their proofs are by hand, plus finite checks that can be
repeated from the printed set \cite{polak2019}.

\paragraph{Data availability.}
The sets, the verifier, the scripts that produce every number here, the preregistrations and the
machine reconciliation of the numbers are archived with the two preprints this paper is assembled from,
\href{https://doi.org/10.5281/zenodo.22979509}{doi:10.5281/zenodo.22979509} and
\href{https://doi.org/10.5281/zenodo.23234251}{doi:10.5281/zenodo.23234251}
\cite{oktiabrev2026a,oktiabrev2026b,shannonrepo}.

\paragraph{Declaration of competing interest.}
The author declares that he has no known competing financial interests or personal relationships that
could have appeared to influence the work reported in this paper.

\paragraph{Funding.}
This research did not receive any specific grant from funding agencies in the public, commercial, or
not-for-profit sectors.

\paragraph{CRediT authorship contribution statement.}
\textbf{Artem Oktiabrev:} conceptualization, methodology, software, validation, formal analysis,
investigation, data curation, writing --- original draft and review.

\paragraph{Declaration of generative AI and AI-assisted technologies in the manuscript preparation
process.}
During the preparation of this work the author used Claude and Claude Code (Anthropic) in order to draft
and revise the text. After using this tool/service the author reviewed and edited the content as needed
and takes full responsibility for the content of the published article. The use of the same tools in the
research itself is described in Section~\ref{sec:repro}.
